\ncdlesson{11}{Folkeskolen mellem faglighed og dannelse}{公立学校在“学科能力”与“全人教养”之间}
{DU3 Modul 5 (PD3) $\rightarrow$ Modul 6 (Studieprøven, B2 $\rightarrow$ C1)｜对标 NCE 4 \& 丹麦教育制度与国民文化内核}
{\textbf{丹麦文化核心不可译词深度辨析}（\textit{dannelse} vs \textit{uddannelse} · \textit{tryghed} vs \textit{sikkerhed} · \textit{fællesskab} vs \textit{samfund}）· 格伦特维（Grundtvig）与哈尔·科赫（Hal Koch）民主对话传统 · 对比并列句式 \textit{Hvor A retter sig mod..., handler B om...}}
{为什么英文和中文都很难用一个词精准对译丹麦语中的 \textit{dannelse}、\textit{tryghed} 与 \textit{fællesskab}？这三组词恰恰构成了丹麦公立基础学校（\textit{folkeskolen}）乃至整个北欧福利社会的精神密码！在 PD3 与 Studieprøven 涉及教育改革、儿童成长或公民价值的命题中，能否精准区分 \textbf{\textit{dannelse}} 与 \textbf{\textit{uddannelse}}、\textbf{\textit{tryghed}} 与 \textbf{\textit{sikkerhed}}、\textbf{\textit{fællesskab}} 与 \textbf{\textit{samfund}}，并援引格伦特维（Grundtvig）与哈尔·科赫（Hal Koch）的民主传统，是考官评判 C1 文化理解深度的核心标尺。}

% ---------------------------------------------------------------------
\begin{foerduleaeser}
带着下面两个问题阅读 Tekst A 和 Tekst B：
\begin{enumerate}
  \item \textit{Hvordan definerer Tekst A forskellen mellem de tre danske begrebspar, og hvordan knytter Grundtvigs og Hal Kochs tanker an til folkeskolens opgave?}（Tekst A 是如何界定这三组丹麦语核心概念之差异的？格伦特维与哈尔·科赫的思想又如何与公立学校的使命相联系？）
  \item \textit{Hvorfor understreger læreren Mette og skolelederen Thomas i Tekst B, at kontrol og høje karakterer ikke i sig selv kan skabe tryghed og demokratisk dannelse?}（在 Tekst B 中，教师梅特和校长托马斯为什么强调管控手段与高分本身并不能自动造就心理安全感与民主教养？）
\end{enumerate}
\end{foerduleaeser}

% ---------------------------------------------------------------------
\ncdsubsec{二、Tekster（文化概念精析长文与教研室对话双语篇）}

\subsection*{Tekst A：\textit{Skolen for livet --- eller skolen for arbejdsmarkedet?}（教育文化与概念辨析长文 · 约 380 词）}

\begin{tipbox}{语篇说明}
本篇深入剖析丹麦公立教育百年争论背后的三组“不可译文化词”。阅读时请重点留意\textbf{\color{NCDNavy}深蓝色标注的三组核心文化概念对}，以及\textbf{\color{NCDRed}深红色标注的概念界定与对比论证结构}。
\end{tipbox}

I årtier har politikere, lærere og forældre diskuteret, hvad den danske folkeskole først og fremmest skal måles på. Bør undervisningen fokusere benhårdt på nationale test, Pisa-placeringer og snæver \da{faglighed}, eller skal skolen i lige så høj grad forme hele og kritiske medborgere? For at forstå kernen i denne debat må man undersøge tre danske nøglebegreber, som næsten er umulige at oversætte med et enkelt ord til andre sprog.

Det første og mest grundlæggende begrebspar er \da{uddannelse} over for \da{dannelse}. \dabold{Hvor} \da{uddannelse} \dabold{retter sig mod} målbare kvalifikationer, eksamener og færdigheder til arbejdsmarkedet, \dabold{handler} \da{dannelse} \dabold{om} menneskelig modning, dømmekraft og evnen til at orientere sig etisk i verden. Allerede i 1800-tallet skelnede N.F.S. Grundtvig skarpt mellem ”skolen for livet” og ”skolen for døden” og fremhævede levende folkelig oplysning frem for mekanisk udenadslære. Efter Anden Verdenskrig videreførte teologen Hal Koch denne arv med sit berømte budskab om, at demokrati ikke blot er en styreform med stemmesedler, men en daglig samtale og livsform, som børn må øve sig i allerede i klasseværelset.

Det andet par, man nødvendigvis må \dabold{skelne mellem}, er \da{sikkerhed} og \da{tryghed}. På mange fremmedsprog dækker ét og samme ord begge betydninger, men i dansk skolekultur er forskellen afgørende. \da{Sikkerhed} drejer sig om ydre og målbare tiltag --- såsom brandalarmer, adgangskontrol eller klare ordensregler i skolegården. \da{Tryghed} derimod betegner en indre oplevelse af tillid, nærvær og visheden om, at man kan række hånden op og begå fejl uden at blive mødt med hån. Uden denne dybe tryghed tør elever sjældent stille de skæve spørgsmål, hvor nysgerrighed og reel læring ellers opstår.

Endelig rummer dansk en vigtig skillelinje mellem \da{samfund} og \da{fællesskab}. \da{Samfundet} udgør den overordnede juridiske og økonomiske ramme af love, skatter og rettigheder, vi alle indgår i som borgere. Et \da{fællesskab} kan derimod ikke vedtages af Folketinget; det opstår gennem frivilligt engagement, fælles erfaringer og gensidig forpligtelse mellem konkrete mennesker --- for eksempel i en blandet folkeskoleklasse, hvor børn fra forskellige sociale lag lærer at lytte til hinanden.

Når skoledebatten spidser til, bliver faglighed og dannelse ofte \dabold{stillet op som modsætninger}. I Grundtvigs og Hal Kochs ånd \dabold{forudsætter de imidlertid hinanden}: En solid uddannelse giver eleverne redskaberne til at handle i samfundet, mens dannelse, tryghed og et levende fællesskab giver dem visdommen og modet til at bruge redskaberne klogt.

\begin{transbox}
\small
\textbf{\color{NCDTeal}[Tekst A 中文精译]}\\
几十年来，政治人物、教师与家长一直在争论：丹麦公立基础学校（Folkeskole）究竟首先应当以什么来衡量？教学应当强硬地聚焦于全国统测、PISA 国际排名和狭义的“学科专业能力”，还是应当在同等程度上塑造完整而具备批判精神的公民？要理解这场辩论的核心，就必须剖析三个几乎无法用单一外语词汇直接翻译的丹麦语关键词。

第一组也是最根本的概念对，是“职业/学业教育（\textit{uddannelse}）”与“全人教养（\textit{dannelse}）”。如果说 \textit{uddannelse} 指向可衡量的资质、考试以及面向劳动力市场的技能，那么 \textit{dannelse} 关乎的则是人的心智成熟、独立判断力以及在世界中进行伦理定位的能力。早在 19 世纪，N.F.S. 格伦特维便尖锐区分了“为了生命的学校”与“死气沉沉的书斋学校”，推崇充满生机的民众启蒙，而非机械的死记硬背。二战结束后，神学家兼思想家哈尔·科赫延续了这一传统，提出了那句著名论断：民主不仅仅是一种凭选票运作的政体，更是一种日常的对话与生活方式，孩子们必须早在教室里就开始练习它。

人们必须严格区分的第二组词，是“外在安全（\textit{sikkerhed}）”与“内在安心感（\textit{tryghed}）”。在许多外语中，这两个意思往往只对应同一个词，但在丹麦校园文化里，二者的差异具有决定性意义。\textit{Sikkerhed} 指的是外在的、技术性的且可衡量的防范措施——例如火警装置、门禁管控或操场上的明确纪律规定。而 \textit{tryghed} 则指一种内在的信任感、真诚在场与确信：即使自己举手答错问题，也不会遭到嘲笑。若缺少这种深层安心感，学生就很少敢于提出那些奇特而发散的问题——而真正的好奇心与深度学习恰恰诞生于此。

最后，丹麦语在“国家社会（\textit{samfund}）”与“价值共同体（\textit{fællesskab}）”之间也保留着重要分界线。\textit{Samfundet} 构成了我们作为公民置身其中的宏观法律与经济框架，由法律、税收与权利组成。相比之下，\textit{fællesskab} 却无法由议会立法直接通过；它产生于具体的人与人之间自愿的投入、共同的经历以及相互的道义承担——例如在一个来自不同社会阶层的混编公立学校班级里，孩子们学会彼此倾听。

当教育辩论趋于白热化时，学科能力与全人教养常常被摆成一对水火不容的对立面。然而，在格伦特维与哈尔·科赫的精神传统中，二者其实互为前提：扎实的学业教育赋予学生进入社会行动的工具，而全人教养、深层安心感与充满生机的共同体，则赋予他们明智运用这些工具的智慧与勇气。
\end{transbox}

\subsection*{Tekst B：\textit{Lærerværelset efter skolebestyrelsesmødet}（校董会后教研室对话 · 约 155 词）}

\textbf{阅读任务（理解与核对）：}\textit{Læs samtalen mellem folkeskolelærer Mette og skoleleder Thomas efter et møde med skolebestyrelsen, og besvar spørgsmålene nedenfor.}

\begin{emailbox}{Skoledebat｜Mette (folkeskolelærer) \& Thomas (skoleleder)}
\textbf{Mette:} Thomas, lagde du mærke til forældrenes uro over vores nye trivselsforløb? Flere spurgte direkte, om vi forsømmer den hårde faglighed, når klassen bruger en hel formiddag på samtaleøvelser og projekter.

\textbf{Thomas:} Ja, det er en klassisk bekymring på mange skoler. Desværre blander mange \da{uddannelse} og \da{dannelse} sammen. Gode karakterer i dansk og matematik er selvfølgelig vigtige, men uden evnen til at lytte, tænke selvstændigt og diskutere uenigheder svigter vi folkeskolens opgave i Hal Kochs forstand.

\textbf{Mette:} Præcis. En far foreslog endda strammere overvågning og faste sanktioner mod uro i timerne. Sådanne tiltag skaber måske ydre \da{sikkerhed}, men de skaber ikke i sig selv den \da{tryghed}, som gør eleverne modige i undervisningen.

\textbf{Thomas:} Helt enig. \da{Samfundet} kan udstede nok så mange læreplaner og love, men det forpligtende \da{fællesskab} må lærere og elever sammen bygge op i klasselokalet hver eneste dag.
\end{emailbox}

\textbf{Tekst B 阅读核对（Rigtigt / Forkert / Kort svar，答案见第六部分）：}
\begin{enumerate}
  \item \textit{Både Mette og Thomas mener, at karakterer i dansk og matematik slet ikke har nogen betydning.} \underline{\hspace{2.2cm}}
  \item \textit{Hvad er problemet med streng overvågning og sanktioner ifølge Mette?} \underline{\hspace{2.2cm}}
  \item \textit{Hvilken opgave kan samfundet ifølge Thomas ikke løse gennem love og læreplaner alene?} \underline{\hspace{2.2cm}}
\end{enumerate}

\begin{tipbox}{本课发音与语调要点（Udtale \& Prosodi）}
\begin{itemize}
  \item \textbf{丹麦文化核心词与人物专名读音}：\textit{dannelse} \ipa{ˈdɑnəlsə}（注意 \textit{nn} 发短元音 \ipa{ɑ}）；\textit{tryghed} \ipa{ˈtˢʁyɡˌheðˀ}（软 \textit{d} 加 \textit{stød}）；\textit{fællesskab} \ipa{ˈfɛləsˌsɡæːˀb}（重音在最前面 \textit{fæl-}）；\textit{N.F.S. Grundtvig} \ipa{ˈɡʁɔnd̥vi}（\textit{dt} 与结尾 \textit{g} 均不发音！）；\textit{Hal Koch} \ipa{hɑl ˈkʰʌɡ}。
  \item \textbf{概念辨析关键动词与名词}：\textit{at skelne} \ipa{ˈsɡɛlnə}；\textit{at forudsætte} \ipa{ˈfɒːuðˌsɛd̥ə}（重音在前缀 \textit{for-}，\textit{d} 读软 \textit{d}）；\textit{dømmekraft} \ipa{ˈdømˀəˌkʰʁɑfd}（\textit{døm-} 带 \textit{stød}）；\textit{udenadslære} \ipa{ˈuːðənɑðsˌlɛːɐ}（\textit{uden-} 与 \textit{-ad-} 均含软 \textit{d}）。
  \item \textbf{对比句式重读节奏}：朗读 \textit{Hvor \textbf{ud}dannelse retter sig mod..., handler \textbf{dan}nelse om...} 时，必须将清晰对比重音分别放在两对反义核心词的首音节上（\textit{\textbf{ud}-} vs \textit{\textbf{dan}-}、\textit{\textbf{sik}-} vs \textit{\textbf{tryg}-}、\textit{\textbf{sam}-} vs \textit{\textbf{fæl}-}）。
\end{itemize}
\end{tipbox}

\begin{transbox}
\small
\textbf{\color{NCDTeal}[Tekst B 中文对照]}\\
\textbf{梅特：}托马斯，你注意到家长们对我们新推行的班级心理福祉课程感到不安了吗？好几位家长直接发问，当全班花整整一上午做对话练习和专题项目时，我们是不是荒废了扎实的学科学习。

\textbf{托马斯：}是的，这是许多学校里都非常经典的顾虑。遗憾的是，很多人把“学业培训（\textit{uddannelse}）”和“全人教养（\textit{dannelse}）”混为一谈了。丹麦语和数学拿到好成绩当然很重要，但如果缺少倾听、独立思考以及讨论分歧的能力，我们就辜负了哈尔·科赫意义上的公立学校使命。

\textbf{梅特：}说得太对了。有位父亲甚至建议在课堂上实行更严格的监控和固定惩戒，好让教室保持安静。这类措施或许能制造出外在的“安全秩序（\textit{sikkerhed}）”，但它们本身并不会自动生发出让学生在课堂上勇敢提问的“内在安心感（\textit{tryghed}）”。

\textbf{托马斯：}完全赞同。“国家社会（\textit{samfundet}）”可以颁布再多的教学大纲和法律条文，但那种彼此承担道义的“共同体（\textit{fællesskab}）”，必须靠老师和学生每一天在教室里共同搭建起来。
\end{transbox}

% ---------------------------------------------------------------------
\ncdsubsec{三、Gloser（核心生词表）}

\begin{tipbox}{词表重点}
本课三大文化词对的“可数性”极易出错：\textit{en dannelse}（全人教养）、\textit{en faglighed}（学科专业性）、\textit{en tryghed}（内在安心感）以及 \textit{en sikkerhed}（外在安全秩序，一般语义下）均属不可数共性名词（\textit{kun ental}）；而 \textit{en uddannelse}（学历/教育项目）可数，复数为 \textit{uddannelser}。另外注意不规则动词：\textit{at forudsætte}（以……为前提）过去式与完成时分别作 \textit{forudsatte, har forudsat}；\textit{at opstå}（出现/产生）完成时用 \dabold{er}。
\end{tipbox}

\begin{glosertable}
名词 & en dannelse & \textit{(dannelsen, kun ental)} & 全人教养，人文陶冶 \\
名词 & en uddannelse & \textit{(uddannelsen, uddannelser, uddannelserne)} & 学业教育，职业培训，学历 \\
名词 & en faglighed & \textit{(fagligheden, kun ental)} & 学科能力，专业水准 \\
名词 & en tryghed & \textit{(trygheden, kun ental)} & 内在安心感，心理安全与信任 \\
名词 & en sikkerhed & \textit{(sikkerheden, kun ental)} & 外在安全，治安防卫保障 \\
名词 & et fællesskab & \textit{(fællesskabet, fællesskaber, fællesskaberne)} & 共同体，凝聚的群体关系 \\
名词 & et samfund & \textit{(samfundet, samfund, samfundene)} & 社会，国家制度共同体 \\
名词 & en dømmekraft & \textit{(dømmekraften, kun ental)} & 判断力，辨别力 \\
名词 & en livsform & \textit{(livsformen, livsformer, livsformerne)} & 生活方式，生命形态 \\
名词 & en udenadslære & \textit{(udenadslæren, kun ental)} & 死记硬背 \\
名词 & en foranstaltning & \textit{(foranstaltningen, -er, foranstaltningerne)} & 措施，防范举措 \\
名词 & et nærvær & \textit{(nærværet, kun ental)} & 真诚在场，陪伴感 \\
\midrule
动词 [har] & at skelne (mellem) & \textit{(skelner, skelnede, har skelnet)} & 区分，辨别 \\
动词 [har] & at forudsætte & \textit{(forudsætter, forudsatte, har forudsat)} & 以……为前提，预设 \\
动词 [har] & at svigte & \textit{(svigter, svigtede, har svigtet)} & 辜负，失职，背弃 \\
动词 [har] & at forsømme & \textit{(forsømmer, forsømte, har forsømt)} & 荒废，疏忽 \\
动词 [er] & at opstå & \textit{(opstår, opstod, er opstået)} & 产生，出现，萌发 \\
\midrule
形容词 & forpligtende & \textit{(forpligtende, 不变)} & 彼此承担责任的，有约束力的 \\
\end{glosertable}

% ---------------------------------------------------------------------
\ncdsubsec{四、Grammatik \& Stilistik（不可译文化词辨析、对比句法与 C1 避坑）}

\begin{grammatikbox}{语法点 1：丹麦文化三大不可译核心词深度辨析矩阵}
\vspace{0.2em}
\renewcommand{\arraystretch}{1.24}
{\small
\begin{tabularx}{\linewidth}{@{} >{\bfseries\raggedright\arraybackslash}p{2.9cm} >{\raggedright\arraybackslash}p{4.6cm} >{\raggedright\arraybackslash}X @{}}
\toprule
\textbf{概念对照对} & \textbf{丹麦文化与制度内核差异} & \textbf{地道语境辨析例句} \\
\midrule
\da{dannelse}\newline vs \da{uddannelse} &
\textit{uddannelse} = 可测量的学历、职业技能与资格；\newline \textit{dannelse} = 人格成熟、判断力、伦理与公民修养（Grundtvig 传统）。 &
\textit{\da{Uddannelse} giver et eksamensbevis, mens \da{dannelse} lærer os at bruge vores viden etisk.} \\
\midrule
\da{tryghed}\newline vs \da{sikkerhed} &
\textit{sikkerhed} = 门锁、监控、警报、制度规则等客观防护；\newline \textit{tryghed} = 信任、接纳、不怕犯错的心理安全感。 &
\textit{Kameraer øger måske den ydre \da{sikkerhed}, men skaber ikke automatisk indre \da{tryghed}.} \\
\midrule
\da{fællesskab}\newline vs \da{samfund} &
\textit{samfund} = 法律、税收与公共福利构成的宏观制度框架；\newline \textit{fællesskab} = 基于参与和互助的情感与责任共同体。 &
\textit{\da{Samfundet} vedtager love, men et forpligtende \da{fællesskab} bygges nedefra i hverdagen.} \\
\bottomrule
\end{tabularx}
}
\end{grammatikbox}

\begin{grammatikbox}{语法点 2：概念界定与对比修辞四大 C1 句型}
\begin{enumerate}
  \item \textbf{对比界定主句铁律：\dabold{Hvor A retter sig mod X, handler B om Y}}：\newline
    注意：\textit{Hvor} 在此处不是空间副词“在哪里”，而是高阶对比连词“如果说 A 指向 X，那么 B 关乎 Y”！前面整个 \textit{Hvor}-从句内部不倒装，但占据主句第 1 位，因此逗号后面主句\textbf{必须 V2 倒装}：\newline
    \ding{51}\ \textit{\da{Hvor} uddannelse retter sig mod eksamener, \dabold{handler dannelse} om dømmekraft.}
  \item \textbf{精准区分：\dabold{at skelne mellem A og B}}：\newline
    后面固定接介词 \dabold{mellem}（不可用 \textit{fra}），对应英文 \textit{distinguish between}：\newline
    \ding{51}\ \textit{Man må nødvendigvis \da{skelne mellem} ydre sikkerhed og indre tryghed.}
  \item \textbf{破除二元对立：\dabold{stille A og B op som modsætninger} vs \dabold{forudsætte hinanden}}：\newline
    先指出舆论的误区（把两者对立起来），再用“互为前提”完成辩证升华：\newline
    \ding{51}\ \textit{Ofte \da{stilles} faglighed og dannelse \da{op som modsætninger}, selvom de i virkeligheden \da{forudsætter hinanden}.}
  \item \textbf{偏好与取舍：\dabold{X frem for Y}}（重 X 轻 Y / 宁要 X 而非 Y）：\newline
    \ding{51}\ \textit{Grundtvig fremhævede levende oplysning \da{frem for} mekanisk udenadslære.}
\end{enumerate}
\end{grammatikbox}

\begin{grammatikbox}{语法点 3：文化概念误用与语序排雷清单}
\begin{enumerate}
  \item \textbf{\ding{172} 切勿用 \textit{sikkerhed} 翻译“班级心理安全感”}：中文的“安全感”若指孩子敢于发言、不怕被霸凌，丹麦语必须用 \da{tryghed}！说 \ding{55}\ \textit{eleverne mangler sikkerhed i klassen} 会让人以为教室门锁坏了或有消防隐患。
  \item \textbf{\ding{173} \textit{Hvor}-对比从句后的 V2 倒装不可漏}：常见错误是受英语 \textit{Whereas A focuses on X, B deals with Y} 影响忘了倒装：\ding{55}\ \textit{Hvor samfundet handler om love, fællesskabet handler om tillid.} $\rightarrow$ \ding{51}\ \textit{Hvor samfundet handler om love, \da{handler fællesskabet} om tillid.}
  \item \textbf{\ding{174} 区分 \textit{at blande A og B sammen} 与 \textit{at skelne mellem A og B}}：“混淆 A 和 B”用可分动词短语 \textit{blande A og B sammen}；“区分 A 和 B”用 \textit{skelne mellem A og B}，切勿写成 \ding{55}\ \textit{skelne A fra B} 或 \ding{55}\ \textit{blande A med B}。
  \item \textbf{\ding{175} 不规则过去分词 \textit{har forudsat}}：\textit{forudsætte} 变化跟 \textit{sætte} 一样（\textit{forudsatte, har forudsat}），注意末尾不加 \textit{-et}，不可误写为 \ding{55}\ \textit{har forudsattet}。
\end{enumerate}
\end{grammatikbox}

% ---------------------------------------------------------------------
\ncdsubsec{五、Mønsterøvelser（句型与修辞操练）}

\textbf{练习 A：丹麦文化核心词与概念对比句式精准填空（每空填入最准确的 1 个丹麦语词）}
\begin{enumerate}
  \item Hvor \underline{\hspace{2.2cm}} giver eleverne konkrete kvalifikationer, former dannelse deres menneskelige dømmekraft.
  \item I debatten om folkeskolen må vi skelne \underline{\hspace{2.2cm}} ydre kontrol og ægte tillid i klassen.
  \item Låste døre og faste ordensregler giver måske teknisk \underline{\hspace{2.2cm}}, men skaber ikke automatisk varme relationer.
  \item Børn tør først række hånden op og begå fejl, når der hersker reel \underline{\hspace{2.2cm}} i klasseværelset. (内在安心感)
  \item Et forpligtende \underline{\hspace{2.2cm}} kan ikke vedtages ved lov, men må gro mellem konkrete mennesker. (共同体)
  \item Grundtvig fremhævede levende samtale frem \underline{\hspace{2.2cm}} mekanisk udenadslære.
  \item Faglighed og dannelse er ikke fjender; i virkeligheden \underline{\hspace{2.2cm}} de hinanden. (forudsætte，现在时)
  \item Hvis skolen kun måler testresultater, \underline{\hspace{2.2cm}} den sin demokratiske opgave i Hal Kochs ånd. (svigte，现在时)
\end{enumerate}

\textbf{练习 B：文化概念辨析纠错与 C1 句法升级（修正误用/语序，或按提示合并）}
\begin{enumerate}
  \item 【文化概念误用纠错】\textit{Når eleverne ikke bliver latterliggjort for fejl, oplever de stor sikkerhed i undervisningen.} $\rightarrow$
  \item 【对比从句 V2 倒装纠错】\textit{Hvor samfundet udgør den juridiske ramme, et fællesskab bygger på frivillig tillid.} $\rightarrow$
  \item 【动词搭配与介词纠错】\textit{Vi bør ikke skelne faglighed fra dannelse, og heller ikke blande dem med hinanden.} $\rightarrow$
  \item 【不规则过去完成时纠错】\textit{Hal Kochs demokratiopfattelse har altid forudsattet en aktiv samtale mellem borgerne.} $\rightarrow$
  \item 【对比句型合并升级】用 \textit{Hvor ... retter sig mod ..., handler ... om ...} 合并：\textit{Uddannelse retter sig mod arbejdsmarkedet. Dannelse handler om etisk dømmekraft.} $\rightarrow$
  \item 【破除二元对立升级】用 \textit{stille ... op som modsætninger} 与 \textit{forudsætte hinanden} 改写：\textit{Politikere siger tit, at test og dannelse ikke passer sammen, men de har brug for hinanden.} $\rightarrow$
\end{enumerate}

\textbf{练习 C：中译丹实战（精准驾驭三大文化词对与对比句型）}
\begin{enumerate}
  \item 在关于公立学校的讨论中，我们必须严格区分学业培训与全人教养。
  \item 如果说学业教育指向具体技能，全人教养则关乎人的道德判断力。
  \item 外在的安全规则本身并不会自动在课堂上创造内在的安心感。
  \item 国家社会可以通过法律，但有凝聚力的共同体却产生于日常相处之中。
  \item 学科能力与全人教养不应被摆成对立面，因为二者其实互为前提。
  \item 在哈尔·科赫的精神中，民主首先是一种必须从小练习的日常对话。
\end{enumerate}

% ---------------------------------------------------------------------
\ncdsubsec{六、PD3 / Studieprøven 实战微任务与参考答案}

\begin{emailbox}{PD3 / Studieprøven 写作微任务（Kultur- og debatkommentar · 200--250+ ord）}
\textbf{Situation}: En politisk tænketank foreslår, at den danske folkeskole fremover skal reducere antallet af kreative fag og klassens time for i stedet at afholde flere nationale test i læsning og matematik samt indføre strengere disciplinære regler.\\[3pt]
\textbf{Opgave}: Skriv en nuanceret debatartikel på \textbf{200--250 ord}, hvor du vurderer forslaget. Brug lektionens tre danske kernebegrebspar (\textit{uddannelse / dannelse, sikkerhed / tryghed, samfund / fællesskab}) aktivt i din argumentation og inddrag kort enten Grundtvigs eller Hal Kochs syn på skolens demokratiske rolle.
\end{emailbox}

\begin{answerbox}
\textbf{Før du læser 参考回答}
\begin{enumerate}
  \item \textit{Tekst A skelner mellem uddannelse (målbare kvalifikationer) og dannelse (modning og dømmekraft), mellem ydre sikkerhed og indre tryghed samt mellem samfundets juridiske ramme og menneskelige fællesskaber. Grundtvig talte om ”skolen for livet”, og Hal Koch så demokrati som en daglig samtale, der skal øves i skolen.}
  \item \textit{De påpeger, at høje karakterer og faglighed ikke kan stå alene uden demokratisk samtaleevne, og at streng overvågning højst skaber ydre sikkerhed, mens reel tryghed og et forpligtende fællesskab må opbygges gennem daglig tillid mellem lærere og elever.}
\end{enumerate}

\textbf{Tekst B 阅读核对答案}
\begin{enumerate}
  \item \da{Forkert}（托马斯明确承认丹麦语和数学的成绩当然也很重要：\textit{Gode karakterer ... er selvfølgelig vigtige}）
  \item \da{Det skaber måske ydre sikkerhed, men ikke i sig selv den tryghed, der gør elever modige.}（严格监控顶多带来外在秩序，却不会自动产生让学生敢于表达的心理安心感）
  \item \da{At bygge det forpligtende fællesskab op i klasselokalet.}（在教室里建构彼此担责的共同体）
\end{enumerate}

\textbf{练习 A 答案}
\begin{enumerate}
  \item \da{uddannelse}（与后文 \textit{dannelse} 构成“学历技能 vs 全人教养”对比）
  \item \da{mellem}（固定介词搭配：\textit{skelne mellem A og B}）
  \item \da{sikkerhed}（对应门禁、规章等客观外在安全措施）
  \item \da{tryghed}（指不怕犯错、彼此信任的内在心理安全感）
  \item \da{fællesskab}（指无法靠立法强制生成、在人与人之间自然生长的共同体）
  \item \da{for}（固定对比取舍短语：\textit{frem for} ……而非……）
  \item \da{forudsætter}（动词 \textit{at forudsætte} 现在时：互为前提）
  \item \da{svigter}（动词 \textit{at svigte} 现在时：辜负、失职）
\end{enumerate}

\textbf{练习 B 答案（参考）}
\begin{enumerate}
  \item \textit{Når eleverne ikke bliver latterliggjort for fejl, oplever de stor \da{tryghed} i undervisningen.}（心理层面安心不怕出错必须用 \textit{tryghed}，不可误用 \textit{sikkerhed}）
  \item \textit{Hvor samfundet udgør den juridiske ramme, \da{bygger et fællesskab} på frivillig tillid.}（\textit{Hvor}-对比从句置于句首时，主句变位动词 \textit{bygger} 须触发 V2 倒装）
  \item \textit{Vi bør \da{skelne mellem} faglighed og dannelse, men heller ikke \da{blande} dem \da{sammen}.}（“区分”固定搭配 \textit{skelne mellem A og B}；“混淆”用 \textit{blande A og B sammen}）
  \item \textit{Hal Kochs demokratiopfattelse har altid \da{forudsat} en aktiv samtale mellem borgerne.}（不规则过去分词作 \textit{forudsat}，不加 \textit{-tet}）
  \item \textit{\da{Hvor} uddannelse retter sig mod arbejdsmarkedet, \da{handler} dannelse om etisk dømmekraft.}（借助 \textit{Hvor ..., handler ...} 构建严密概念界定句并正确倒装）
  \item \textit{Politikere \da{stiller} ofte test og dannelse \da{op som modsætninger}, selvom de faktisk \da{forudsætter hinanden}.}（以固定短语替代口语化表述，突出“互为前提”的辩证力）
\end{enumerate}

\textbf{练习 C 答案}
\begin{enumerate}
  \item \textit{I debatten om folkeskolen må vi skelne klart mellem uddannelse og dannelse.}
  \item \textit{Hvor uddannelse retter sig mod konkrete færdigheder, handler dannelse om etisk dømmekraft.}
  \item \textit{Ydre regler om sikkerhed skaber ikke i sig selv indre tryghed i klassen.}
  \item \textit{Samfundet vedtager love, men et stærkt fællesskab opstår gennem dagligt samvær.}
  \item \textit{Faglighed og dannelse bør ikke stilles op som modsætninger, da de forudsætter hinanden.}
  \item \textit{I Hal Kochs ånd er demokrati først og fremmest en samtale, som skal øves.}
\end{enumerate}

\textbf{PD3 / Studieprøven 高分范文（约 235 词，不含标题）}\\
\textit{\textbf{En stærk folkeskole kræver både faglighed og levende dannelse}}

\textit{Forslaget om flere nationale test og strengere disciplin på bekostning af kreative fag og klassens time rammer skævt, fordi det bygger på en falsk modsætning. Ganske vist er gode læse- og regnefærdigheder uundværlige for alle unge, men en folkeskole reduceret til ren testtræning overser selve kernen i den danske skoletradition.}

\textit{Her er det helt afgørende at skelne klart mellem \da{uddannelse} og \da{dannelse}. \dabold{Hvor uddannelse retter sig mod} målbare kundskaber og fremtidig beskæftigelse, \dabold{handler dannelse om} dømmekraft, nysgerrighed og evnen til at indgå konstruktivt i demokratisk uenighed. Allerede i 1945 mindede Hal Koch os om, at demokrati først og fremmest er en livsform og en daglig samtale, som børn må øve sig i gennem undervisningen --- ikke en mekanisk eksamensdisciplin, man kan lære ved udenadslære.}

\textit{Samtidig forveksler forslaget ydre \da{sikkerhed} med indre \da{tryghed}. Strenge disciplinære regler og øget kontrol kan måske dæmpe uroen kortvarigt, men reel læring og mod opstår kun, når eleverne føler sig så trygge i klassens \da{fællesskab}, at de tør begå fejl og stille skæve spørgsmål uden frygt for hån. \da{Samfundet} kan udstede nok så mange test, læreplaner og sanktioner, men det kan aldrig vedtage engagement og gensidig tillid ved lov.}

\textit{Derfor bør faglighed og dannelse ikke stilles op som modsætninger i folkeskolen. De forudsætter tværtimod hinanden: Kun når børn trives i trygge, forpligtende fællesskaber med plads til samtale og kreativ udfoldelse, får de også det mentale overskud til at løfte deres faglige uddannelse.}

\small\textbf{名师点评}：正文（不含标题）约 235 词，文化纵深感极强！文章开篇否定智库提案的虚假二元对立，随后自然流畅地展开三组丹麦不可译文化词的精确辨析：\textit{uddannelse} vs \textit{dannelse}（援引哈尔·科赫 1945 年民主对话观）、\textit{sikkerhed} vs \textit{tryghed}（揭示纪律控制与班级安全感之别）、\textit{samfund} vs \textit{fællesskab}（点明立法与共同体的边界），最后收束于“互为前提”，是 Studieprøven 教育文化题的满分样板！
\end{answerbox}